\documentclass[12pt,a4paper]{article}

%% ============================================================
%% PACKAGES — all standard, available in any TeX Live / MiKTeX
%% ============================================================
\usepackage[top=2.5cm,bottom=2.5cm,left=2.8cm,right=2.8cm]{geometry}
\usepackage{amsmath,amssymb}
\usepackage{xcolor}
\usepackage{booktabs}
\usepackage{longtable}
\usepackage{array}
\usepackage{enumitem}
\usepackage{setspace}
\usepackage{fancyhdr}
\usepackage{titlesec}
\usepackage{tcolorbox}
\usepackage{hyperref}

%% Fix headheight
\setlength{\headheight}{15pt}

%% ============================================================
%% COLORS
%% ============================================================
\definecolor{darkblue}{RGB}{10,45,100}
\definecolor{midblue}{RGB}{30,90,160}
\definecolor{lightblue}{RGB}{220,235,255}
\definecolor{darkgreen}{RGB}{20,90,50}
\definecolor{lightgreen}{RGB}{220,245,225}
\definecolor{darkorange}{RGB}{160,70,10}
\definecolor{lightorange}{RGB}{255,240,220}
\definecolor{darkred}{RGB}{140,20,20}
\definecolor{lightred}{RGB}{255,225,225}
\definecolor{lightgray}{RGB}{245,245,248}
\definecolor{midgray}{RGB}{150,150,160}
\definecolor{darkgray}{RGB}{60,60,70}

%% ============================================================
%% TCOLORBOX STYLES
%% ============================================================
\tcbuselibrary{skins,breakable}

\tcbset{
  milestonestyle/.style={
    colback=lightred,
    colframe=darkred,
    boxrule=1.2pt,
    arc=4pt,
    left=8pt,right=8pt,top=6pt,bottom=6pt,
    breakable
  },
  deliverablestyle/.style={
    colback=lightgreen,
    colframe=darkgreen,
    boxrule=1.2pt,
    arc=4pt,
    left=8pt,right=8pt,top=6pt,bottom=6pt,
    breakable
  },
  warningstyle/.style={
    colback=lightorange,
    colframe=darkorange,
    boxrule=1.2pt,
    arc=4pt,
    left=8pt,right=8pt,top=6pt,bottom=6pt,
    breakable
  },
  infostyle/.style={
    colback=lightblue,
    colframe=midblue,
    boxrule=1.2pt,
    arc=4pt,
    left=8pt,right=8pt,top=6pt,bottom=6pt,
    breakable
  },
  daystyle/.style={
    colback=lightgray,
    colframe=midgray,
    boxrule=0.6pt,
    arc=3pt,
    left=8pt,right=8pt,top=5pt,bottom=5pt,
    breakable
  }
}

%% ============================================================
%% SECTION FORMATTING
%% ============================================================
\titleformat{\section}
  {\Large\bfseries\color{darkblue}}
  {\thesection.}{0.8em}{}
  [\vspace{-0.3em}{\color{midblue}\rule{\linewidth}{1.5pt}}]

\titleformat{\subsection}
  {\large\bfseries\color{midblue}}
  {\thesubsection.}{0.6em}{}

%% ============================================================
%% HEADER AND FOOTER
%% ============================================================
\pagestyle{fancy}
\fancyhf{}
\fancyhead[L]{\small\color{darkgray}\textit{TL-NN Research Plan --- I.~Sedka}}
\fancyhead[R]{\small\color{darkgray}\textit{Laghouat, Algeria --- 2025}}
\fancyfoot[C]{\small\color{midgray}\thepage}
\renewcommand{\headrulewidth}{0.4pt}

%% ============================================================
%% CUSTOM COMMANDS
%% ============================================================
\newcommand{\hrs}[1]{\textit{\color{darkgray}(#1~hours)}}

\newcommand{\weekheading}[2]{%
  \vspace{0.8em}
  \noindent
  {\large\bfseries\color{midblue}--- Week~#1 ---
  \hfill\normalsize\normalfont\color{darkgray}#2}
  \vspace{0.2em}
  \\{\color{midblue}\rule{\linewidth}{0.6pt}}
  \vspace{0.3em}
}

\newcommand{\dayheading}[2]{%
  \vspace{0.3em}
  \noindent{\bfseries\color{darkblue}Day~#1}%
  \hfill{\small\color{darkgray}\textit{#2}}%
  \\[-0.4em]{\color{midgray}\rule{\linewidth}{0.4pt}}%
  \vspace{0.1em}
}

%% ============================================================
\begin{document}
\onehalfspacing
%% ============================================================

%% ------------------------------------------------------------
%% TITLE PAGE
%% ------------------------------------------------------------
\begin{titlepage}
\begin{center}

\vspace*{1.5cm}
{\color{darkblue}\rule{\linewidth}{2pt}}
\vspace{0.5cm}

{\Huge\bfseries\color{darkblue} Three-Month Research Plan}\\[0.6em]
{\Large\color{midblue} Writing and Publishing the TL-NN Paper}

\vspace{0.5cm}
{\color{midblue}\rule{\linewidth}{1pt}}
\vspace{1.2cm}

\begin{tcolorbox}[infostyle]
\begin{center}
{\large\bfseries\color{darkblue}
A Hybrid Transform--Linearize--Neural Network Method\\[4pt]
for Nonlinear Fredholm Integral Equations on Large Intervals}
\end{center}
\end{tcolorbox}

\vspace{1.5cm}

\begin{tabular}{rl}
\textbf{\color{darkgray}Researcher:}    & Ilyes SEDKA \\[4pt]
\textbf{\color{darkgray}Position:}      & Assistant Professor in Applied Mathematics \\[4pt]
\textbf{\color{darkgray}Institution:}   & Amar Telidji University, Laghouat, Algeria \\[4pt]
\textbf{\color{darkgray}Co-author:}     & Ammar KHELLAF, Natl.\ Polytechnic College Constantine \\[4pt]
\textbf{\color{darkgray}Target:}        & Applied Mathematics and Computation (Q2, Elsevier) \\[4pt]
\textbf{\color{darkgray}Duration:}      & 12 weeks $\times$ 3 hours/day $= 252$ working hours \\[4pt]
\textbf{\color{darkgray}Date:}          & \today \\
\end{tabular}

\vspace{1.8cm}

\begin{tcolorbox}[warningstyle]
\noindent\textbf{How to use this document.}
Read the relevant section every morning before starting work.
Mark completed tasks.
If you miss a day, do not double up ---
continue from where you stopped.
Every section ends with a concrete deliverable.
If you cannot produce the deliverable, you have not
finished the section.
\end{tcolorbox}

\vfill
{\small\color{midgray}Private document --- not for distribution.}

\end{center}
\end{titlepage}

%% ------------------------------------------------------------
\newpage
\tableofcontents
\newpage

%% ============================================================
\section{Strategic Overview}
%% ============================================================

\subsection{The Idea in One Paragraph}

The TL-NN method solves nonlinear Fredholm integral equations
of the second kind on large intervals $[0,\tau]$, $\tau \gg 1$.
It is built directly from the TLD algorithm of \cite{sedka2025}
by preserving the Transform phase (domain decomposition into $N$
subintervals) and the Linearize phase (Newton--Kantorovich
iteration) exactly as published, while replacing only the
Nystr\"{o}m discretization phase with a small feedforward neural
network trained at each Newton iteration to approximate the
Fr\'{e}chet derivative of the integral operator.
The convergence framework of TLD is inherited entirely.
The paper applies TL-NN to the three benchmark examples from
\cite{sedka2025} and compares results directly with TLD.

\subsection{Why This Paper Is the Right First Step}

\begin{itemize}[leftmargin=1.8em,itemsep=4pt]
  \item The mathematical foundation already exists in
    \cite{sedka2025}. You extend a proven framework,
    not build a new one from scratch.
  \item The numerical baseline already exists. Tables~1--3
    of \cite{sedka2025} provide the TLD values. You add
    TL-NN columns next to each table.
  \item The scope is calibrated precisely for three months
    of part-time work at three hours per day.
  \item It opens the research direction honestly. The
    conclusion points toward TLD-XPINN and a full
    convergence theorem for neural network training.
\end{itemize}

\subsection{Target Journal}

\begin{center}
\begin{tabular}{>{\bfseries\color{darkgray}}p{4cm} p{8.5cm}}
\toprule
Journal       & Applied Mathematics and Computation \\[2pt]
Publisher     & Elsevier \\[2pt]
Quartile      & Q2 (Scimago) \\[2pt]
ISSN          & 0096-3003 \\[2pt]
Impact factor & 3.5 (2024) \\[2pt]
Review time   & 6--10 weeks typical \\[2pt]
Template      & \texttt{elsarticle.cls} \\[2pt]
Target length & 15--18 pages \\
\bottomrule
\end{tabular}
\end{center}

\vspace{0.5em}
\begin{tcolorbox}[infostyle]
\noindent\textbf{Why this journal.}
Applied Mathematics and Computation publishes exactly
this type of paper: a new numerical method with convergence
analysis and numerical examples. Your reference \cite{khellaf2022}
is already published there, which confirms your work fits
the scope.
\end{tcolorbox}

\subsection{Time Budget}

\begin{center}
\begin{tabular}{lccc}
\toprule
\textbf{Phase} & \textbf{Weeks} & \textbf{Days} & \textbf{Hours} \\
\midrule
Month~1 --- Learning and Implementation & 1--4  & 28 & 84 \\
Month~2 --- Writing and Revision        & 5--8  & 28 & 84 \\
Month~3 --- Submission and Next Steps   & 9--12 & 28 & 84 \\
\midrule
\textbf{Total} & \textbf{12} & \textbf{84} & \textbf{252} \\
\bottomrule
\end{tabular}
\end{center}

\begin{tcolorbox}[warningstyle]
\noindent\textbf{Non-negotiable rule on time.}
Three hours per day means three hours of deep, uninterrupted
work. Phone in another room. No email, no social media.
A physical timer running.
\end{tcolorbox}

\subsection{Paper Structure}

\begin{center}
\begin{tabular}{clcc}
\toprule
\textbf{Sec.} & \textbf{Title} & \textbf{Pages} & \textbf{Written in} \\
\midrule
1 & Introduction           & 2  & Week~5 \\
2 & Preliminaries          & 3  & Weeks~1 and~3 \\
3 & The TL-NN Algorithm    & 4  & Week~3 \\
4 & Convergence Discussion & 2  & Week~3 \\
5 & Numerical Results      & 4  & Week~4 \\
6 & Conclusion             & 1  & Week~3 \\
  & References             & 2  & Week~8 \\
\midrule
  & \textbf{Total}         & \textbf{18} & \\
\bottomrule
\end{tabular}
\end{center}

%% ============================================================
\newpage
\section{Month~1: Learning and Implementation}
%% ============================================================

\begin{tcolorbox}[infostyle]
\noindent\textbf{Month~1 Objective.}
By the end of Week~4:
(1)~the PINNs paper understood and Subsection~2.2 written
in \LaTeX;
(2)~a working Python implementation of TL-NN for all three
examples from \cite{sedka2025};
(3)~complete error tables ready to insert into the paper.
\end{tcolorbox}

\weekheading{1}{Reading --- PINNs and Neural Network Basics}

\noindent\textbf{Goal:}
Understand the PINNs paper sufficiently to write
Subsection~2.2 of the paper. No implementation this week.

\vspace{0.5em}

\begin{tcolorbox}[daystyle]
\dayheading{1}{3 hours}
\begin{itemize}[leftmargin=1.8em,itemsep=3pt]
  \item[\color{midblue}$\bullet$]
    \hrs{1} Watch ``But what is a neural network?''
    by 3Blue1Brown (YouTube, 19~min). After: write one
    sentence explaining what a neural network computes,
    using only functional analysis language.
  \item[\color{midblue}$\bullet$]
    \hrs{1} Watch ``Gradient descent, how neural networks
    learn'' by 3Blue1Brown (21~min). After: write one
    sentence connecting gradient descent to the
    Newton--Kantorovich iteration in \cite{sedka2025}.
  \item[\color{midblue}$\bullet$]
    \hrs{1} Download the PINNs paper (Raissi, Perdikaris,
    Karniadakis 2019, arXiv:1711.10561). Read only the
    abstract and Section~1. Write one paragraph: what
    problem do they solve and what is their method?
\end{itemize}
\noindent\textbf{End of day:} One page of handwritten notes.
\end{tcolorbox}

\begin{tcolorbox}[daystyle]
\dayheading{2}{3 hours}
\begin{itemize}[leftmargin=1.8em,itemsep=3pt]
  \item[\color{midblue}$\bullet$]
    \hrs{1.5} Read Sections~2 and~3 of the PINNs paper.
    Recognize their loss function as a residual norm.
    Build a two-column translation table:
    \textit{their term} vs.\ \textit{your term}.
  \item[\color{midblue}$\bullet$]
    \hrs{1.5} Read Sections~4 and~5. Compare their error
    magnitudes with your TLD error tables from
    \cite{sedka2025}. Write three honest observations.
\end{itemize}
\noindent\textbf{End of day:}
Translation table (at least 8 rows) + three observations.
\end{tcolorbox}

\begin{tcolorbox}[daystyle]
\dayheading{3}{3 hours}
\begin{itemize}[leftmargin=1.8em,itemsep=3pt]
  \item[\color{midblue}$\bullet$]
    \hrs{3} Write a one-page summary of the PINNs paper in
    functional analysis language. No ML jargon. Then write
    one paragraph answering: what does the PINNs paper not
    prove that your TLD paper does prove? This paragraph
    becomes the core argument of your paper's introduction.
\end{itemize}
\noindent\textbf{End of day:}
One-page summary + one paragraph on the convergence gap.
\end{tcolorbox}

\begin{tcolorbox}[daystyle]
\dayheading{4}{3 hours}
\begin{itemize}[leftmargin=1.8em,itemsep=3pt]
  \item[\color{midblue}$\bullet$]
    \hrs{1.5} Read the DeepONet paper (Lu et al.\ 2021,
    arXiv:1910.03193): abstract, Section~1, and
    Definition~2.1 only. Write Definition~2.1 next to
    your operator $K_l$ from \cite{sedka2025}. Are they
    defining the same object? Write your honest answer.
  \item[\color{midblue}$\bullet$]
    \hrs{1.5} Read their discussion section. Write every
    limitation they acknowledge. For each one, state how
    your TLD framework addresses it.
\end{itemize}
\noindent\textbf{End of day:}
Operator comparison + limitation list.
\end{tcolorbox}

\begin{tcolorbox}[daystyle]
\dayheading{5}{3 hours}
\begin{itemize}[leftmargin=1.8em,itemsep=3pt]
  \item[\color{midblue}$\bullet$]
    \hrs{3} Write Subsection~2.2 (Physics-Informed Neural
    Networks) in \LaTeX. Open \texttt{TL\_NN\_paper.tex}
    and write the text. Target: one page. Explain PINNs
    in mathematical language, state why direct PINN
    application on large intervals fails (spectral bias),
    and explain why TL-NN uses a simpler learning task.
\end{itemize}
\noindent\textbf{End of day:}
Subsection~2.2 complete and compiled without errors.
\end{tcolorbox}

\begin{tcolorbox}[daystyle]
\dayheading{6}{3 hours}
\begin{itemize}[leftmargin=1.8em,itemsep=3pt]
  \item[\color{midblue}$\bullet$]
    \hrs{1.5} Install Python environment. Required:
    \texttt{numpy}, \texttt{matplotlib}, \texttt{scipy},
    \texttt{torch}. Command:
    \texttt{pip install numpy matplotlib scipy torch}.
    Verify each package imports without error.
  \item[\color{midblue}$\bullet$]
    \hrs{1.5} Run your first PyTorch program: define a
    2-hidden-layer network with 20 neurons and $\tanh$
    activation; generate 100 training points
    $(x_i, \sin x_i)$ for $x_i \in [0, 2\pi]$;
    train for 1000 epochs; plot result vs.\ true function.
\end{itemize}
\noindent\textbf{End of day:}
Python and PyTorch installed.
$\sin(x)$ approximation plot saved.
\end{tcolorbox}

\begin{tcolorbox}[daystyle]
\dayheading{7}{3 hours}
\begin{itemize}[leftmargin=1.8em,itemsep=3pt]
  \item[\color{midblue}$\bullet$]
    \hrs{1} Build one master correspondence table:
    left column --- ML concepts from PINNs and DeepONet;
    right column --- equivalent concept in your TLD work.
  \item[\color{midblue}$\bullet$]
    \hrs{2} Write the first rough draft of Section~1
    (Introduction) in \LaTeX. For each paragraph
    placeholder, write at least three sentences.
    Mark unclear parts with \texttt{\%~TODO}.
    Do not try to make it good --- make it exist.
\end{itemize}
\noindent\textbf{End of day:}
Correspondence table + rough introduction draft in \LaTeX.
\end{tcolorbox}

\begin{tcolorbox}[milestonestyle]
\noindent\textbf{\color{darkred}Week~1 Milestone.}
PINNs paper understood. Subsection~2.2 written and compiled.
PyTorch installed and tested. Rough introduction exists
in \LaTeX.
\end{tcolorbox}

\weekheading{2}{Implementation --- First Working Code}

\noindent\textbf{Goal:}
Write Python code implementing TL-NN and produce
complete error tables for all three examples.

\vspace{0.5em}

\begin{tcolorbox}[daystyle]
\dayheading{8}{3 hours}
\begin{itemize}[leftmargin=1.8em,itemsep=3pt]
  \item[\color{midblue}$\bullet$]
    \hrs{1} Study your TLD MATLAB code. Identify the exact
    lines that build matrix $H_n^{(k)}$. Write their
    line numbers and purpose in a comment file.
    These are the only lines you will replace.
  \item[\color{midblue}$\bullet$]
    \hrs{2} Translate TLD from MATLAB to Python for
    Example~1 (TLD only, no neural network yet). Verify
    you reproduce the error from Table~1 of
    \cite{sedka2025} for $n = 5$, $\tau = 20$
    within 10\% tolerance.
\end{itemize}
\noindent\textbf{End of day:}
Python TLD code running and verified for Example~1.
\end{tcolorbox}

\begin{tcolorbox}[daystyle]
\dayheading{9}{3 hours}
\begin{itemize}[leftmargin=1.8em,itemsep=3pt]
  \item[\color{midblue}$\bullet$]
    \hrs{3} Write the neural network module
    \texttt{kernel\_net.py}. Input: training pairs
    $\{(t_i, t_j, \partial_{\psi_p}\kappa_l(t_i,t_j))\}$
    and hyperparameters. Output: trained network evaluable
    at any $(t,s)$. Train on the kernel derivative of
    Example~1 at $k = 0$ (zero iterate). Verify loss
    decreases. Plot and save the training loss curve.
\end{itemize}
\noindent\textbf{End of day:}
\texttt{kernel\_net.py} module + training loss plot saved.
\end{tcolorbox}

\begin{tcolorbox}[daystyle]
\dayheading{10}{3 hours}
\begin{itemize}[leftmargin=1.8em,itemsep=3pt]
  \item[\color{midblue}$\bullet$]
    \hrs{1.5} Integrate \texttt{kernel\_net.py} into the
    TLD iteration. Replace the Nystr\"{o}m evaluation of
    $H_n^{(k)}$ with network evaluation.
    Everything else stays identical.
  \item[\color{midblue}$\bullet$]
    \hrs{1.5} Run full TL-NN for Example~1, $n = 5$,
    $\tau = 20$. Record: Newton iteration count, error
    $E_{n,N}$, and computation time. Compare with TLD.
\end{itemize}
\noindent\textbf{End of day:}
First TL-NN result: one number in the table.
\end{tcolorbox}

\begin{tcolorbox}[daystyle]
\dayheading{11}{3 hours}
\begin{itemize}[leftmargin=1.8em,itemsep=3pt]
  \item[\color{midblue}$\bullet$]
    \hrs{3} Run TL-NN for Example~1 across all
    $n \in \{5, 10, 30, 50, 100\}$ for $\tau = 20$.
    Fill the first column of Table~1 in
    \texttt{TL\_NN\_paper.tex}.
    Do not spend more than 30 minutes debugging
    any single run.
\end{itemize}
\noindent\textbf{End of day:}
First column of Table~1 filled in \LaTeX.
\end{tcolorbox}

\begin{tcolorbox}[daystyle]
\dayheading{12}{3 hours}
\begin{itemize}[leftmargin=1.8em,itemsep=3pt]
  \item[\color{midblue}$\bullet$]
    \hrs{1.5} Run TL-NN for Example~1, $\tau = 50$,
    all $n$ values. Fill second column of Table~1.
  \item[\color{midblue}$\bullet$]
    \hrs{1.5} Run TL-NN for Example~1, $\tau = 80$,
    all $n$ values. Fill third column. Table~1 complete.
\end{itemize}
\noindent\textbf{End of day:} Table~1 complete in \LaTeX.
\end{tcolorbox}

\begin{tcolorbox}[daystyle]
\dayheading{13}{3 hours}
\begin{itemize}[leftmargin=1.8em,itemsep=3pt]
  \item[\color{midblue}$\bullet$]
    \hrs{3} Run TL-NN for Example~2
    ($\kappa = \cos(3s)\,e^{(t+3s)/30}\psi^3(s)$)
    for all $n$ and $\tau \in \{10, 50, 80\}$.
    Fill Table~2 in \texttt{TL\_NN\_paper.tex}.
\end{itemize}
\noindent\textbf{End of day:} Table~2 complete in \LaTeX.
\end{tcolorbox}

\begin{tcolorbox}[daystyle]
\dayheading{14}{3 hours}
\begin{itemize}[leftmargin=1.8em,itemsep=3pt]
  \item[\color{midblue}$\bullet$]
    \hrs{2} Run TL-NN for Example~3
    ($\kappa = \tfrac{1}{4}te^{-ts}
    [2(1 - \cos\pi s) + \psi^2(s)]$)
    for all $n$ and $\tau \in \{20, 50, 100\}$.
    Fill Table~3 in \texttt{TL\_NN\_paper.tex}.
  \item[\color{midblue}$\bullet$]
    \hrs{1} Verify all three tables: TLD columns must
    match published values from \cite{sedka2025}.
    Write a one-paragraph summary of what the results
    show so far.
\end{itemize}
\noindent\textbf{End of day:}
All three tables complete and verified.
\end{tcolorbox}

\begin{tcolorbox}[milestonestyle]
\noindent\textbf{\color{darkred}Week~2 Milestone.}
All three numerical tables complete in \LaTeX.
TL-NN produces output for all three examples at all
values of $n$ and $\tau$.
\end{tcolorbox}

\weekheading{3}{Writing --- Sections 3, 4, and 6}

\noindent\textbf{Goal:}
Write Sections~3, 4, and~6 of the paper, produce all
figures, and send the draft to Khellaf.

\vspace{0.5em}

\begin{tcolorbox}[daystyle]
\dayheading{15}{3 hours}
\begin{itemize}[leftmargin=1.8em,itemsep=3pt]
  \item[\color{midblue}$\bullet$]
    \hrs{1} Analyze the three tables honestly. For each
    example and each $\tau$: does TL-NN outperform,
    match, or underperform TLD? Write bullet-point answers.
    This becomes the first draft of Subsection~5.4.
  \item[\color{midblue}$\bullet$]
    \hrs{2} Retrain with larger $N_{\max}$ or smaller
    $\varepsilon_{\mathrm{train}}$ on any runs where
    TL-NN significantly underperforms. Update tables.
\end{itemize}
\noindent\textbf{End of day:}
Analysis notes + updated tables if needed.
\end{tcolorbox}

\begin{tcolorbox}[daystyle]
\dayheading{16}{3 hours}
\begin{itemize}[leftmargin=1.8em,itemsep=3pt]
  \item[\color{midblue}$\bullet$]
    \hrs{3} Produce Figure~1: for Example~1, $n = 30$,
    $\tau = 20$, plot TLD approximate solution, TL-NN
    approximate solution, and exact solution
    $\psi_{\mathrm{ext}}(t)$ on the same axes.
    Three distinguishable curves, legend, axis labels.
    Save as \texttt{fig1.pdf} at 300~dpi.
\end{itemize}
\noindent\textbf{End of day:}
Figure~1 saved and imported into \LaTeX.
\end{tcolorbox}

\begin{tcolorbox}[daystyle]
\dayheading{17}{3 hours}
\begin{itemize}[leftmargin=1.8em,itemsep=3pt]
  \item[\color{midblue}$\bullet$]
    \hrs{1.5} Produce Figure~2: for Example~1, $n = 30$,
    $\tau = 20$, plot $\log_{10}(E_n^{(k)})$ vs.\
    Newton iteration $k$ for TLD and TL-NN.
    Confirms linear convergence. Save as \texttt{fig2.pdf}.
  \item[\color{midblue}$\bullet$]
    \hrs{1.5} Produce Figure~3: training loss
    $\mathcal{L}(\theta)$ vs.\ epoch number for Example~2,
    $n = 10$, $\tau = 50$. Save as \texttt{fig3.pdf}.
\end{itemize}
\noindent\textbf{End of day:}
Figures~2 and~3 saved and imported into \LaTeX.
\end{tcolorbox}

\begin{tcolorbox}[daystyle]
\dayheading{18}{3 hours}
\begin{itemize}[leftmargin=1.8em,itemsep=3pt]
  \item[\color{midblue}$\bullet$]
    \hrs{3} Write Section~3 (The TL-NN Algorithm) in
    \LaTeX. The equation skeleton already exists in
    \texttt{TL\_NN\_paper.tex}. Write the connecting
    prose: motivating paragraph before each equation and
    interpreting paragraph after it. Every step in
    Algorithm~1 must reference its equation number.
    Target: four pages, no \texttt{\%~TODO} markers.
\end{itemize}
\noindent\textbf{End of day:}
Section~3 complete, four pages, compiled.
\end{tcolorbox}

\begin{tcolorbox}[daystyle]
\dayheading{19}{3 hours}
\begin{itemize}[leftmargin=1.8em,itemsep=3pt]
  \item[\color{midblue}$\bullet$]
    \hrs{3} Write Section~4 (Convergence Discussion) in
    \LaTeX. State Proposition~1 and Theorem~1 with
    proofs mirroring those of \cite{sedka2025}, replacing
    $\delta_{n,N}$ with $\varepsilon_{\mathrm{NN}}$.
    Write the remark comparing the convergence rates of
    TLD and TL-NN. Target: two pages.
\end{itemize}
\noindent\textbf{End of day:}
Section~4 complete, two pages, compiled.
\end{tcolorbox}

\begin{tcolorbox}[daystyle]
\dayheading{20}{3 hours}
\begin{itemize}[leftmargin=1.8em,itemsep=3pt]
  \item[\color{midblue}$\bullet$]
    \hrs{3} Write Subsection~2.1 (The TLD Framework) in
    \LaTeX. Adapt text from \cite{sedka2025} directly
    (legitimate self-citation). Ensure every equation
    uses the same notation as Section~3.
    Target: three pages.
\end{itemize}
\noindent\textbf{End of day:}
Subsection~2.1 complete, three pages, compiled.
\end{tcolorbox}

\begin{tcolorbox}[daystyle]
\dayheading{21}{3 hours}
\begin{itemize}[leftmargin=1.8em,itemsep=3pt]
  \item[\color{midblue}$\bullet$]
    \hrs{1.5} Write Section~6 (Conclusion): (1)~what was
    done; (2)~what the results show; (3)~four specific
    future directions: TLD-XPINN, weakly singular kernels,
    systems of equations, full convergence theorem for
    neural network training.
  \item[\color{midblue}$\bullet$]
    \hrs{1.5} Send Ammar Khellaf an email with the current
    PDF. Request review of Sections~3 and~4 for
    mathematical correctness. Deadline: two weeks.
\end{itemize}
\noindent\textbf{End of day:}
Section~6 complete + email sent to Khellaf.
\end{tcolorbox}

\begin{tcolorbox}[milestonestyle]
\noindent\textbf{\color{darkred}Week~3 Milestone.}
Sections~2.1, 3, 4, and~6 written in \LaTeX.
All three figures produced and imported.
Email sent to Khellaf with current draft.
\end{tcolorbox}

\weekheading{4}{Completing the Rough Draft}

\noindent\textbf{Goal:}
Write Section~5, complete the abstract, verify all content,
and produce the Month~1 complete rough draft.

\vspace{0.5em}

\begin{tcolorbox}[daystyle]
\dayheading{22}{3 hours}
\begin{itemize}[leftmargin=1.8em,itemsep=3pt]
  \item[\color{midblue}$\bullet$]
    \hrs{3} Write Subsections~5.1, 5.2, and~5.3. For
    each example: one paragraph on the equation and exact
    solution; one paragraph on experimental parameters;
    reference to the table and figure; two to three
    observation sentences below the table.
\end{itemize}
\noindent\textbf{End of day:} Subsections~5.1--5.3 complete.
\end{tcolorbox}

\begin{tcolorbox}[daystyle]
\dayheading{23}{3 hours}
\begin{itemize}[leftmargin=1.8em,itemsep=3pt]
  \item[\color{midblue}$\bullet$]
    \hrs{3} Write Subsection~5.4 (Discussion and
    Comparison). Address five points: (1)~quantitative
    comparison across all examples; (2)~effect of
    increasing $\tau$; (3)~training cost vs.\ quadrature
    cost; (4)~when TL-NN outperforms TLD and when it does
    not; (5)~connection to Section~4 convergence rates.
    Target: one page.
\end{itemize}
\noindent\textbf{End of day:} Subsection~5.4 complete.
\end{tcolorbox}

\begin{tcolorbox}[daystyle]
\dayheading{24}{3 hours}
\begin{itemize}[leftmargin=1.8em,itemsep=3pt]
  \item[\color{midblue}$\bullet$]
    \hrs{3} Compile the complete paper for the first time
    as a single document. Read from beginning to end
    without editing. Mark every problem with a
    \texttt{\%~FIX} comment. Categorize: mathematical /
    writing / notation.
\end{itemize}
\noindent\textbf{End of day:}
First complete compiled draft with categorized markers.
\end{tcolorbox}

\begin{tcolorbox}[daystyle]
\dayheading{25}{3 hours}
\begin{itemize}[leftmargin=1.8em,itemsep=3pt]
  \item[\color{midblue}$\bullet$]
    \hrs{3} Fix all mathematical \texttt{\%~FIX} markers.
    Every proof step must follow from something before it.
    Every equation referenced correctly.
    Every assumption stated before first use.
\end{itemize}
\noindent\textbf{End of day:}
Zero mathematical \texttt{\%~FIX} markers remaining.
\end{tcolorbox}

\begin{tcolorbox}[daystyle]
\dayheading{26}{3 hours}
\begin{itemize}[leftmargin=1.8em,itemsep=3pt]
  \item[\color{midblue}$\bullet$]
    \hrs{3} Write the abstract. Structure: (1)~context;
    (2)~the gap; (3)~what TL-NN does; (4)~main results;
    (5)~theoretical contribution. Write five 150-word
    drafts. Choose the best. Insert into
    \texttt{TL\_NN\_paper.tex}.
\end{itemize}
\noindent\textbf{End of day:}
Abstract complete, 150 words, compiled.
\end{tcolorbox}

\begin{tcolorbox}[daystyle]
\dayheading{27}{3 hours}
\begin{itemize}[leftmargin=1.8em,itemsep=3pt]
  \item[\color{midblue}$\bullet$]
    \hrs{1.5} Verify the bibliography. Every cited reference
    has a \texttt{bibitem}. Every \texttt{bibitem} cited
    at least once. All DOI links present.
    Target: at least 19 entries.
  \item[\color{midblue}$\bullet$]
    \hrs{1.5} Check notation consistency. Every symbol
    defined before first use. Pay special attention to
    $n$ vs.\ $N$, $m$ vs.\ $M$, and $k$ as Newton index.
\end{itemize}
\noindent\textbf{End of day:}
Bibliography verified + notation consistent.
\end{tcolorbox}

\begin{tcolorbox}[daystyle]
\dayheading{28}{3 hours}
\begin{itemize}[leftmargin=1.8em,itemsep=3pt]
  \item[\color{midblue}$\bullet$]
    \hrs{3} Compile the final Month~1 draft. Count pages
    (target: 15--20). Zero compilation errors. Zero
    mathematical \texttt{\%~FIX} markers. Save as
    \texttt{TL\_NN\_draft\_month1.pdf}.
\end{itemize}
\noindent\textbf{End of day:}
\texttt{TL\_NN\_draft\_month1.pdf} ready for Month~2.
\end{tcolorbox}

\begin{tcolorbox}[deliverablestyle]
\noindent\textbf{\color{darkgreen}Month~1 Deliverable.}
Complete rough draft: all six sections, all three tables,
all three figures, bibliography, abstract. Compiled as a
clean 15--20 page PDF. Zero compilation errors. Zero
mathematical errors. File: \texttt{TL\_NN\_draft\_month1.pdf}.
\end{tcolorbox}

%% ============================================================
\newpage
\section{Month~2: Writing and Revision}
%% ============================================================

\begin{tcolorbox}[infostyle]
\noindent\textbf{Month~2 Objective.}
By the end of Week~8: a polished, submission-ready draft
reviewed by Khellaf, revised twice, fully compliant with
journal requirements, and submitted to Applied Mathematics
and Computation.
\end{tcolorbox}

\weekheading{5}{First Revision --- Writing Quality}

\begin{tcolorbox}[daystyle]
\dayheading{29}{3 hours}
\begin{itemize}[leftmargin=1.8em,itemsep=3pt]
  \item[\color{midblue}$\bullet$]
    \hrs{3} Read Section~1 (Introduction) aloud. Fix
    every awkward sentence. Rule for every paragraph:
    first sentence states the topic; every sentence
    advances the argument; last sentence transitions.
\end{itemize}
\noindent\textbf{End of day:}
Section~1 revised: clear and tight.
\end{tcolorbox}

\begin{tcolorbox}[daystyle]
\dayheading{30}{3 hours}
\begin{itemize}[leftmargin=1.8em,itemsep=3pt]
  \item[\color{midblue}$\bullet$]
    \hrs{1.5} Apply same revision to Section~2
    (Preliminaries).
  \item[\color{midblue}$\bullet$]
    \hrs{1.5} Apply same revision to Section~6
    (Conclusion).
\end{itemize}
\noindent\textbf{End of day:} Sections~2 and~6 revised.
\end{tcolorbox}

\begin{tcolorbox}[daystyle]
\dayheading{31}{3 hours}
\begin{itemize}[leftmargin=1.8em,itemsep=3pt]
  \item[\color{midblue}$\bullet$]
    \hrs{3} Revise Section~3 (TL-NN Algorithm). Algorithm~1
    must match the text exactly: every variable defined in
    the text, every step referenced by equation number.
\end{itemize}
\noindent\textbf{End of day:}
Section~3 consistent with Algorithm~1.
\end{tcolorbox}

\begin{tcolorbox}[daystyle]
\dayheading{32}{3 hours}
\begin{itemize}[leftmargin=1.8em,itemsep=3pt]
  \item[\color{midblue}$\bullet$]
    \hrs{3} Revise Section~4 (Convergence). Test: can you
    defend every proof step to an expert referee? If not:
    add the justification or honestly flag it as future work.
\end{itemize}
\noindent\textbf{End of day:}
Section~4: every step justified or honestly flagged.
\end{tcolorbox}

\begin{tcolorbox}[daystyle]
\dayheading{33}{3 hours}
\begin{itemize}[leftmargin=1.8em,itemsep=3pt]
  \item[\color{midblue}$\bullet$]
    \hrs{3} Revise Section~5 (Numerical Results). Verify
    every table number against raw Python output.
    One wrong digit is grounds for rejection --- verify
    twice. Check every figure: axis labels, legend, font
    size, file format (PDF).
\end{itemize}
\noindent\textbf{End of day:}
Every number confirmed. Every figure checked.
\end{tcolorbox}

\begin{tcolorbox}[daystyle]
\dayheading{34}{3 hours}
\begin{itemize}[leftmargin=1.8em,itemsep=3pt]
  \item[\color{midblue}$\bullet$]
    \hrs{3} If Khellaf has responded: categorize each
    comment as (A)~mathematical fix, (B)~writing
    suggestion, or (C)~clarification needed. Address all
    (A) and (B) comments. If no response: send a reminder.
\end{itemize}
\noindent\textbf{End of day:}
All co-author comments addressed in \LaTeX.
\end{tcolorbox}

\begin{tcolorbox}[daystyle]
\dayheading{35}{3 hours}
\begin{itemize}[leftmargin=1.8em,itemsep=3pt]
  \item[\color{midblue}$\bullet$]
    \hrs{1.5} Compile Revision~1. Save as
    \texttt{TL\_NN\_revision1.pdf}. Send to Khellaf
    for final mathematical confirmation. Deadline: one week.
  \item[\color{midblue}$\bullet$]
    \hrs{1.5} Read the journal author guidelines at the
    Elsevier website. Note: required sections, reference
    style, figure format, ethical statements.
    Write a submission checklist.
\end{itemize}
\noindent\textbf{End of day:}
Revision~1 sent to Khellaf + submission checklist written.
\end{tcolorbox}

\begin{tcolorbox}[milestonestyle]
\noindent\textbf{\color{darkred}Week~5 Milestone.}
First complete revision finalized. All sections clear and
tight. Every table number verified. Revision~1 sent to
Khellaf. Submission checklist ready.
\end{tcolorbox}

\weekheading{6}{Formatting and Journal Compliance}

\begin{tcolorbox}[daystyle]
\dayheading{36}{3 hours}
\begin{itemize}[leftmargin=1.8em,itemsep=3pt]
  \item[\color{midblue}$\bullet$]
    \hrs{1} Download the Elsevier \texttt{elsarticle}
    \LaTeX{} template from the Elsevier website.
  \item[\color{midblue}$\bullet$]
    \hrs{2} Convert \texttt{TL\_NN\_paper.tex} to use
    \texttt{elsarticle.cls}. Reformat author, affiliation,
    and abstract. Compile and verify.
\end{itemize}
\noindent\textbf{End of day:}
Paper compiled with \texttt{elsarticle.cls}.
\end{tcolorbox}

\begin{tcolorbox}[daystyle]
\dayheading{37}{3 hours}
\begin{itemize}[leftmargin=1.8em,itemsep=3pt]
  \item[\color{midblue}$\bullet$]
    \hrs{1.5} Convert bibliography to
    \texttt{elsarticle-num.bst}. Full journal names,
    DOI links, complete page numbers.
  \item[\color{midblue}$\bullet$]
    \hrs{1.5} Save each figure as a separate file:
    \texttt{fig1.pdf}, \texttt{fig2.pdf}, \texttt{fig3.pdf}.
    Verify each is at least $7 \times 5$~cm at 300~dpi.
\end{itemize}
\noindent\textbf{End of day:}
Bibliography and figures compliant with journal requirements.
\end{tcolorbox}

\begin{tcolorbox}[daystyle]
\dayheading{38}{3 hours}
\begin{itemize}[leftmargin=1.8em,itemsep=3pt]
  \item[\color{midblue}$\bullet$]
    \hrs{3} Add all required Elsevier statements:
    (1)~Declaration of Competing Interest;
    (2)~CRediT author statement;
    (3)~Acknowledgements (thank DGRSDT);
    (4)~Data availability statement:
    ``Code and data available from the corresponding
    author upon reasonable request.''
\end{itemize}
\noindent\textbf{End of day:}
All required statements added to \LaTeX.
\end{tcolorbox}

\begin{tcolorbox}[daystyle]
\dayheading{39}{3 hours}
\begin{itemize}[leftmargin=1.8em,itemsep=3pt]
  \item[\color{midblue}$\bullet$]
    \hrs{3} Read the paper as a referee would. Ask:
    (1)~Is the contribution clearly stated?
    (2)~Is Algorithm~1 implementable from the paper alone?
    (3)~Are the experiments sufficient?
    (4)~Is the convergence discussion honest about what
    is proved? Fix every ``no'' answer.
\end{itemize}
\noindent\textbf{End of day:}
Referee self-review complete. All issues fixed.
\end{tcolorbox}

\begin{tcolorbox}[daystyle]
\dayheading{40}{3 hours}
\begin{itemize}[leftmargin=1.8em,itemsep=3pt]
  \item[\color{midblue}$\bullet$]
    \hrs{3} Write the cover letter. One page, three
    paragraphs. \textbf{Paragraph~1}: what TL-NN
    contributes and why it is original. \textbf{Paragraph~2}:
    why this journal is appropriate. \textbf{Paragraph~3}:
    originality and author approval statement.
    State facts only. Do not oversell.
\end{itemize}
\noindent\textbf{End of day:}
Cover letter complete as \texttt{cover\_letter.pdf}.
\end{tcolorbox}

\begin{tcolorbox}[daystyle]
\dayheading{41--42}{3 hours per day}
\begin{itemize}[leftmargin=1.8em,itemsep=3pt]
  \item[\color{midblue}$\bullet$]
    Day~41: Address the second round of Khellaf's comments.
    If not received: strengthen Section~4 by adding a
    remark connecting $\varepsilon_{\mathrm{NN}}$ to
    established approximation results
    (cite Cybenko 1989).
  \item[\color{midblue}$\bullet$]
    Day~42: Compile the paper and verify the submission
    checklist line by line. Create an Elsevier Editorial
    Manager account. Prepare the file list: manuscript,
    \LaTeX{} source, figure files, cover letter.
\end{itemize}
\noindent\textbf{End of Days~41--42:}
Checklist complete. Account ready.
\end{tcolorbox}

\begin{tcolorbox}[milestonestyle]
\noindent\textbf{\color{darkred}Week~6 Milestone.}
Paper formatted with \texttt{elsarticle}. All required
statements included. Cover letter complete.
Editorial Manager account ready.
\end{tcolorbox}

\weekheading{7}{Deep Revision}

\begin{tcolorbox}[daystyle]
\dayheading{43--45}{3 hours per day}
\begin{itemize}[leftmargin=1.8em,itemsep=3pt]
  \item[\color{midblue}$\bullet$]
    Day~43: Read Sections~1 and~2 aloud one final time.
    Cut every sentence that does not earn its place.
  \item[\color{midblue}$\bullet$]
    Day~44: Read Algorithm~1 and Section~3. Test: can you
    implement TL-NN from this section alone without your
    Python code? If not: add what is missing.
  \item[\color{midblue}$\bullet$]
    Day~45: Read Sections~4--6. Every claim in the
    conclusion must be supported by a result in Sections~4
    or~5.
\end{itemize}
\end{tcolorbox}

\begin{tcolorbox}[daystyle]
\dayheading{46}{3 hours}
\begin{itemize}[leftmargin=1.8em,itemsep=3pt]
  \item[\color{midblue}$\bullet$]
    \hrs{3} Final notation check, equation by equation:
    every symbol defined before first use; every equation
    numbered; every equation referenced at least once.
\end{itemize}
\noindent\textbf{End of day:} Zero notation problems.
\end{tcolorbox}

\begin{tcolorbox}[daystyle]
\dayheading{47}{3 hours}
\begin{itemize}[leftmargin=1.8em,itemsep=3pt]
  \item[\color{midblue}$\bullet$]
    \hrs{3} Final language check: subject-verb agreement;
    correct articles (a, an, the); consistent tense
    (present for mathematics, past for experiments);
    no informal contractions. When unsure: simplify.
\end{itemize}
\noindent\textbf{End of day:} Language verified throughout.
\end{tcolorbox}

\begin{tcolorbox}[daystyle]
\dayheading{48}{3 hours}
\begin{itemize}[leftmargin=1.8em,itemsep=3pt]
  \item[\color{midblue}$\bullet$]
    \hrs{1.5} Compile the final PDF. Save as
    \texttt{TL\_NN\_final.pdf}. Pages: 15--18. Zero
    warnings. Zero broken references.
  \item[\color{midblue}$\bullet$]
    \hrs{1.5} Send \texttt{TL\_NN\_final.pdf} to Khellaf:
    ``This is the final version. Please confirm approval
    for submission within one week.''
\end{itemize}
\noindent\textbf{End of day:}
\texttt{TL\_NN\_final.pdf} produced and sent.
\end{tcolorbox}

\weekheading{8}{Submission}

\begin{tcolorbox}[daystyle]
\dayheading{49--51}{3 hours per day}
\begin{itemize}[leftmargin=1.8em,itemsep=3pt]
  \item[\color{midblue}$\bullet$]
    If Khellaf requests corrections: implement now.
    Mathematical corrections only at this stage.
  \item[\color{midblue}$\bullet$]
    If no corrections: prepare \texttt{tlnn\_code.py},
    a clean documented Python script producing all
    tables and figures from the paper.
\end{itemize}
\end{tcolorbox}

\begin{tcolorbox}[daystyle]
\dayheading{52}{3 hours}
\begin{itemize}[leftmargin=1.8em,itemsep=3pt]
  \item[\color{midblue}$\bullet$]
    \hrs{3} Submit through Elsevier Editorial Manager.
    Upload: (1)~cover letter; (2)~main manuscript;
    (3)~\LaTeX{} source; (4)~figure files.
    Article type: Research Article. Suggest three
    referees. Complete all required fields. Submit.
\end{itemize}
\noindent\textbf{End of day:}
\textbf{\color{darkred}Paper submitted.}
Record the submission number and date.
\end{tcolorbox}

\begin{tcolorbox}[daystyle]
\dayheading{53--56}{3 hours per day}
\begin{itemize}[leftmargin=1.8em,itemsep=3pt]
  \item[\color{midblue}$\bullet$]
    Begin reading for the second paper (TLD-XPINN).
    Read: Jagtap, Kharazmi, Karniadakis (2020),
    Conservative physics-informed neural networks on
    discrete domains, CMAME 365, 113028.
    Apply the same reading method as Week~1: translation
    table, honest summary, connection to your TLD work.
\end{itemize}
\noindent\textbf{End of Days~53--56:}
XPINN paper read. First notes toward second paper exist.
\end{tcolorbox}

\begin{tcolorbox}[deliverablestyle]
\noindent\textbf{\color{darkgreen}Month~2 Deliverable.}
Paper submitted to Applied Mathematics and Computation.
Submission number recorded.
Reading for the second paper (TLD-XPINN) begun.
\end{tcolorbox}

%% ============================================================
\newpage
\section{Month~3: Waiting, Responding, and Next Steps}
%% ============================================================

\begin{tcolorbox}[infostyle]
\noindent\textbf{Month~3 Objective.}
Manage the review process professionally, respond to
referees if required, and use the waiting period
productively by advancing the second paper (TLD-XPINN).
\end{tcolorbox}

\subsection{Understanding the Review Timeline}

\begin{center}
\begin{tabular}{p{3.8cm} p{3cm} p{5.5cm}}
\toprule
\textbf{Stage} & \textbf{Duration} & \textbf{Action} \\
\midrule
Editorial check  & 1--2 weeks  & None --- wait \\
Under review     & 4--8 weeks  & Advance next paper \\
First decision   & 6--10 weeks & Respond within 6 weeks \\
Revision         & If required & --- \\
Final decision   & 2--4 weeks  & None --- wait \\
\bottomrule
\end{tabular}
\end{center}

\vspace{0.5em}
\begin{tcolorbox}[infostyle]
\noindent\textbf{Possible first decisions.}
\textbf{Accept} (under 5\%).
\textbf{Major revision} (50--60\%): the paper has merit
but needs significant work. This is a positive outcome,
not a rejection.
\textbf{Minor revision} (25--30\%): small corrections,
straightforward.
\textbf{Reject} (15--20\%): revise and submit to the next
journal on your list.
The correct response to major revision is not
disappointment --- it is focus on addressing every comment
precisely.
\end{tcolorbox}

\weekheading{9--10}{Productive Waiting}

\begin{tcolorbox}[daystyle]
\dayheading{57--70}{3 hours per day}
\begin{itemize}[leftmargin=1.8em,itemsep=3pt]
  \item[\color{midblue}$\bullet$]
    Days~57--59: Complete reading the XPINN paper.
    Write honest summary and correspondence table.
  \item[\color{midblue}$\bullet$]
    Days~60--62: Design the TLD-XPINN algorithm. Write
    the mathematical formulation of the coupled loss
    function for the Fredholm integral coupling.
  \item[\color{midblue}$\bullet$]
    Days~63--65: Create \texttt{TLD\_XPINN\_skeleton.tex}
    with section titles and one sentence per section.
  \item[\color{midblue}$\bullet$]
    Days~66--68: Implement a first version of TLD-XPINN
    for Example~1 in Python with $N = 2$ subintervals.
  \item[\color{midblue}$\bullet$]
    Days~69--70: Write the technical memo for TLD-XPINN:
    what we claim, proof strategy, and open problems.
\end{itemize}
\noindent\textbf{End of Weeks~9--10:}
TLD-XPINN has a skeleton, a Python prototype, and a memo.
\end{tcolorbox}

\weekheading{11--12}{Responding to Referees}

\noindent When the referee report arrives, follow this
protocol exactly.

\vspace{0.5em}

\begin{tcolorbox}[daystyle]
\dayheading{Upon receipt}{Referee response protocol}
\begin{itemize}[leftmargin=1.8em,itemsep=3pt]
  \item[\color{midblue}$\bullet$]
    \textbf{Step~1.} Read once, do not react.
    Wait 24 hours before reading again.
  \item[\color{midblue}$\bullet$]
    \textbf{Step~2.} Categorize every comment:
    (A)~correct and easy;
    (B)~correct and requires work;
    (C)~needs clarification;
    (D)~incorrect --- polite technical rebuttal.
  \item[\color{midblue}$\bullet$]
    \textbf{Step~3.} Fix all (A) comments in Week~1 of
    revision. Plan and complete (B) in Week~2.
  \item[\color{midblue}$\bullet$]
    \textbf{Step~4.} Write the response letter. For each
    comment: quote it exactly, state your response, and
    give the change location (page, line, equation).
    Never argue without technical justification.
  \item[\color{midblue}$\bullet$]
    \textbf{Step~5.} Resubmit within 6 weeks.
    Most journals allow 8 --- resubmit in 6 to leave
    buffer.
\end{itemize}
\end{tcolorbox}

\begin{tcolorbox}[daystyle]
\dayheading{71--84}{3 hours per day}
\begin{itemize}[leftmargin=1.8em,itemsep=3pt]
  \item[\color{midblue}$\bullet$]
    If reports received: follow the protocol above.
    Pause TLD-XPINN work during revision.
  \item[\color{midblue}$\bullet$]
    If reports not yet received: continue TLD-XPINN.
    Target: complete the Python implementation for all
    three examples and produce the first error table
    comparing TLD, TL-NN, and TLD-XPINN on Example~1.
\end{itemize}
\end{tcolorbox}

\begin{tcolorbox}[deliverablestyle]
\noindent\textbf{\color{darkgreen}Month~3 Deliverable.}
Either: revision submitted to the journal (one step from
acceptance), or TLD-XPINN advanced to the point where a
second submission is feasible within six months.
In both cases: the first paper is submitted and the second
paper has meaningfully begun.
\end{tcolorbox}

%% ============================================================
\newpage
\section{Non-Negotiable Rules}
%% ============================================================

\begin{enumerate}[leftmargin=2em,itemsep=8pt,
                  label=\textbf{R\arabic*.}]

  \item \textbf{Three hours means three hours of deep,
    uninterrupted work.} Set a physical timer. Phone in
    another room. No email, no social media during work.

  \item \textbf{Every day ends with a written deliverable.}
    Something in a file, not in your head. If you cannot
    produce the stated deliverable, you have not finished
    the day.

  \item \textbf{If you miss a day: continue, do not
    double up.} A plan executed at 90\% is a success.
    A plan that collapses trying to make up for missed
    days is a failure.

  \item \textbf{Write in \LaTeX{} from Day~18 onward.}
    Not Word. \LaTeX{} forces precision and produces a
    submission-ready document directly.

  \item \textbf{No new reading after Day~14.} Relevant
    papers found later: write the citation and read it
    after submission.

  \item \textbf{Verify every table number twice against
    raw code output before submission.} One wrong digit
    is grounds for rejection.

  \item \textbf{Do not change the paper scope after
    Day~28.} Extensions go into the future work section.

  \item \textbf{Update Khellaf every two weeks.}
    A co-author kept informed responds quickly when you
    need them.

  \item \textbf{Submit even if the paper is not perfect.}
    Referees find imperfections --- that is their job.
    A submitted imperfect paper can become accepted.
    A perfect unsubmitted paper contributes nothing.

  \item \textbf{Keep a log of every date and decision.}
    Submission date, submission number, first decision,
    revision date, final decision.
    This record is required for your academic file.

\end{enumerate}

%% ============================================================
\newpage
\section{Summary Timeline}
%% ============================================================

\begin{center}
\begin{longtable}{p{1.4cm} p{1.4cm} p{2.8cm} p{7cm}}
\toprule
\textbf{Week} & \textbf{Month} & \textbf{Phase}
  & \textbf{Key deliverable} \\
\midrule
\endhead
\midrule
\multicolumn{4}{r}{\small\textit{continued\ldots}}\\
\endfoot
\bottomrule
\endlastfoot
1  & 1 & Reading
   & PINNs paper understood. Subsection~2.2 written.
     PyTorch installed. \\[4pt]
2  & 1 & Implementation
   & TL-NN code running.
     All three error tables filled. \\[4pt]
3  & 1 & Writing
   & Sections~2.1, 3, 4, 6 written.
     Three figures produced. Email to Khellaf. \\[4pt]
4  & 1 & Draft
   & \texttt{TL\_NN\_draft\_month1.pdf}: 15--20 pages,
     zero errors. \\[4pt]
\midrule
5  & 2 & Revision~1
   & All sections revised. Co-author comments addressed.
     Checklist written. \\[4pt]
6  & 2 & Formatting
   & \texttt{elsarticle} applied. Statements included.
     Cover letter written. \\[4pt]
7  & 2 & Deep revision
   & \texttt{TL\_NN\_final.pdf} sent for approval. \\[4pt]
8  & 2 & \textbf{Submission}
   & \textbf{Paper submitted.}
     XPINN reading begun. \\[4pt]
\midrule
9--10 & 3 & Waiting
   & TLD-XPINN skeleton, prototype, and memo. \\[4pt]
11--12 & 3 & Response
   & Referee revision submitted \textit{or}
     TLD-XPINN implementation advanced. \\
\end{longtable}
\end{center}

\begin{tcolorbox}[warningstyle]
\noindent\textbf{Honest expectation.}
If this plan is executed at 80\% or better, by the end of
Month~3 you will have:
(1)~one submitted paper in a Q2 journal;
(2)~a working Python implementation of TL-NN;
(3)~a meaningful start on the second paper (TLD-XPINN);
and (4)~the technical vocabulary to work at the intersection
of numerical analysis and scientific machine learning.
This is the foundation. Every large journey begins with
one submitted paper.
\end{tcolorbox}

%% ============================================================
\newpage
\begin{thebibliography}{99}

\bibitem{sedka2025}
I.~Sedka, A.~Khellaf,
Toward efficient numerical solutions of non-linear integral
equations with TLD algorithm,
\textit{Izvestiya: Mathematics} \textbf{89}(2) (2025)
341--358.

\bibitem{sedka2023}
I.~Sedka, A.~Khellaf, M.Z.~Aissaoui,
New algorithm to solve nonlinear functional equations
applying linearization then double discretization scheme
(L.D.D),
\textit{Kuwait Journal of Science} \textbf{50}(2) (2023)
65--74.

\bibitem{sedka2022a}
I.~Sedka, S.~Lemita, M.Z.~Aissaoui,
Linearization-Discretization process to solve systems of
nonlinear Fredholm integral equations in an
infinite-dimensional context,
\textit{Advances in the Theory of Nonlinear Analysis and
its Applications} \textbf{6}(4) (2022) 547--564.

\bibitem{sedka2022b}
I.~Sedka, A.~Khellaf, S.~Lemita, M.Z.~Aissaoui,
New Algorithm on Linearization-Discretization Solving
Systems of Nonlinear Integro-Differential Fredholm
Equations,
\textit{Boletim da Sociedade Paranaense de
Matem\'{a}tica} \textbf{42} (2024) 1--17.

\bibitem{khellaf2022}
A.~Khellaf, M.Z.~Aissaoui,
New theoretical conditions for solving functional
nonlinear equations by linearization then discretization,
\textit{International Journal of Nonlinear Analysis and
Applications} \textbf{13}(1) (2022) 2857--2869.

\bibitem{lemita2020}
S.~Lemita, H.~Guebbai, I.~Sedka, M.Z.~Aissaoui,
New method for the numerical solution of the Fredholm
linear integral equation on a large interval,
\textit{Russian Universities Reports. Mathematics}
\textbf{25}(132) (2020) 387--400.

\bibitem{raissi2019}
M.~Raissi, P.~Perdikaris, G.E.~Karniadakis,
Physics-informed neural networks: A deep learning
framework for solving forward and inverse problems,
\textit{Journal of Computational Physics}
\textbf{378} (2019) 686--707.

\bibitem{lu2021}
L.~Lu, P.~Jin, G.~Pang, Z.~Zhang, G.E.~Karniadakis,
Learning nonlinear operators via DeepONet,
\textit{Nature Machine Intelligence} \textbf{3} (2021)
218--229.

\bibitem{jagtap2020}
A.D.~Jagtap, E.~Kharazmi, G.E.~Karniadakis,
Conservative physics-informed neural networks on discrete
domains for conservation laws,
\textit{Computer Methods in Applied Mechanics and
Engineering} \textbf{365} (2020) 113028.

\bibitem{atkinson1997}
K.E.~Atkinson,
\textit{The Numerical Solution of Integral Equations of
the Second Kind},
Cambridge University Press, Cambridge, 1997.

\bibitem{grammont2013}
L.~Grammont,
Nonlinear integral equation of the second kind:
a new version of Nystr\"{o}m method,
\textit{Numerical Functional Analysis and Optimization}
\textbf{34}(5) (2013) 496--515.

\bibitem{kingma2015}
D.P.~Kingma, J.L.~Ba,
Adam: A method for stochastic optimization,
in: \textit{Proceedings of ICLR 2015},
arXiv:1412.6980, 2015.

\bibitem{cybenko1989}
G.~Cybenko,
Approximation by superpositions of a sigmoidal function,
\textit{Mathematics of Control, Signals and Systems}
\textbf{2}(4) (1989) 303--314.

\end{thebibliography}

\end{document}
